\documentclass[11pt]{article}
\usepackage{fontspec}
\setmainfont{Times New Roman}
\setsansfont{Arial}
\setmonofont{Consolas}
\usepackage{amsmath,amssymb}
\usepackage{geometry}
\usepackage{microtype}
\geometry{a4paper,margin=3cm}
\setlength{\parindent}{1.25em}
\setlength{\parskip}{0pt}
\emergencystretch=30em
\allowdisplaybreaks
\raggedbottom
\providecommand{\p}{\partial}
\providecommand{\modu}[1]{\;(\operatorname{mod} #1)}

\begin{document}

% Unit: Paper 16. Direct Interslavic Latin translation from the German final-audited source slice.

\section*{16. K razkladam v rědy v teoriji form}
\begin{center}
\emph{Mathematische Annalen} 81 (1920), s. 25--30
\end{center}

Pod «razkladom v rědy v teoriji form», kako jest znano, razuměje se, s jednoj strany, obobčenje razklada Klebš--Gordana: to jest razklad formy, ktora vsebuje dva binarne rědy, po stepenjah determinanta tyh rędov, pri čem koeficienty sut polary form samo od jednogo rěda proměnnyh, ili, čto jest to samo, isčezajut pri priměnjenju \(\Omega\)-procesa. To se obobčaje na formy od libovoljno mnogih serij po \(n\) proměnnyh. S drugoj strany, kako obobčenje normalnyh form, ktore nastupajut pri «kompleksnyh koordinatah» \(t_{ik}\), razuměje se razklad po normalnyh formah za formy, ktore jednovrěmenno vsebujejut veličiny različnyh stupnjev \(t_i,t_{ik},t_{ikl},\ldots\). Ale tuty drugi vid razklada možno dobiti iz prvogo pomoću invariantnyh procesov diferenciranja, kako to prěd vsego izveli Mertens i Derujts.\footnote{F. Mertens, \emph{Über invariante Gebilde quaternärer Formen} (Sitzber. d. Ak. d. Wiss. Wien 98, Abt. IIa (1889)) za kvaternarny slučaj. Obče normalne formy nahodet se u J. Deruyts: \emph{Sur la réduction des fonctions invariantes dans le système des variables géométriques}, \emph{Bulletins de l'Ac. R. de Belgique} 24 (1892). Bez znanja tutoj raboty, v točnoj opori na Mertensa, dala jesm obči razklad po normalnyh formah v \S~7 raboty: \emph{Zur Invariantentheorie der Formen von \(n\) Variabeln}, J. f. M. 139 (1910).}

Različne razklady prvogo vida vse nastajut, kako to kratko načrtano na primer v \emph{Math. Ann.} 77, cit. město III, kako slědstva razklada formy od \(n\) rędov po \(n\) proměnnyh po stepenjah determinanta tyh rędov, pri čem koeficienty sut polary form samo od \((n-1)\) rędov, ktore same izvodet se iz izhodnoj formy tvorjenjem polarov i \(\Omega\)-procesom.\footnote{Sravni bibliografične ukazanja na citovanom městu. Od toga časa pojavila se takože, na formalnoj metodi osnovana, rabota: Schouten, \emph{Over reeksontwikkelingen van algebraische vormen met verschillende rijen van variabelen van verschillenden graad}, \emph{Verslag Ak. Amsterdam} 27 (1919), s. 1481.} Prvotny razklad Klebš--Gordana \((n=2)\) šel ješče jeden krok dalje: specialne polarny procesy, ktore tamo nastupajut, podane sut eksplicitno, čak s čislovymi koeficientami.

To, čto teper hoču dodati k naznačenoj zamětki, gdě razklad v rědy jest pojmany više pojmovno, kako problem linearnyh snopev form, jest eksplicitno predstavjenje do čislovyh koeficientov, ktoro se dobiva specializacijeju izslědovanj E. Fišera\footnote{E. Fischer, \emph{Über die Differentiationsprozesse der Algebra}, J. f. M. 148 (1917).} o občih modulnyh sistemah.

K tomuto vključanju vodi tolkovanje normalnoj formy v \(t_{ik}\). Imennno, ako zaměniti «kompleksne koordinaty» \(t_{ik}\) «linijnymi koordinatami»
\[
        p_{ik}=(x_i y_k-x_k y_i),
\]
togda zaradi identičnyh relacij medžu \(p_{ik}\) predstavjenje formy \(F(p_{ik})\) uže ne jest jednoznačno. Jednoznačnosť možno dostignuti, ako se zahtěvaje, že medžu koeficientami \(F\) važe te same linearne relacije, kako medžu produktami stepenev \(p_{ik}\), ktore nastupajut v \(F\).\footnote{Sravni Clebsch, \emph{Über eine Fundamentalaufgabe der Invariantentheorie}, \S~4, Abh. d. K. G. d. Wissensch. zu Göttingen 17 (1872).} Zato za tutu normalnu formu \(\Phi\) imajemo:
\begin{equation}
\begin{aligned}
        F(p_{ik})&=\Phi(p_{ik}) &&[\text{identično v }x,y],\quad\text{ili}\notag\\
        F(t_{ik})&\equiv\Phi(t_{ik}) &&\modu{M},
\end{aligned}
\tag{1}
\end{equation}
gdě \(M\) označaje modul vseh identičnyh relacij medžu \(p_{ik}\), to jest sostavjen jest iz vsih form \(G(t_{ik})\), za ktore \(G(p_{ik})\) identično v \(x,y\) isčezaje. \(\Phi(t_{ik})\) jest «normalna forma» v kompleksnyh koordinatah, dokolě bazisne funkcije modula \(M\) stajut kvadratičnymi formami v \(t_{ik}\). Iteracijeju postupka vzhodno k faktoram tyh bazisnyh funkcij nastaje polny razklad v rědy. Odpovědajuče razsmatranja važe pri jednovrěmennom nastupanju \(t_i,t_{ik},t_{ikl},\ldots\).

Na analogičny hod mysljenja možno svesti razklad po polarah. Ako, na primer, v formě \(F(x,y,z)\) zaměniti tri ternarne rědy \(x=(x_1,x_2,x_3),y,z\) takymi rędami \(\xi,\eta,\zeta\), ktore linearno sostavljajut se samo iz dvuh iz njih, na primer \(\xi=x,\eta=y;\zeta=\lambda_1x+\lambda_2y\), togda predstavjenje \(F(\xi,\eta,\zeta)\) znovu ne jest jednoznačno. Možno znovu predpisati za koeficienty te same linearne relacije, kako za produkty stepenev \(\xi,\eta,\zeta\), ktore nastupajut v \(F\). Poněže tude modul \(M\) vseh identičnyh relacij imaje kako bazisnu funkciju determinant \(\Delta=(xyz)\) trěh rędov (cit. město III, pomočna lema a)), tude vměsto [1] nastupaje:
\begin{equation}
\begin{aligned}
        F(\xi,\eta,\zeta)&=\Phi(\xi,\eta,\zeta)&&[\text{identično v }x,y],\quad\text{ili}\notag\\
        F(x,y,z)&\equiv\Phi(x,y,z)&&\modu{\Delta}.
\end{aligned}
\tag{2}
\end{equation}
Pokazuje se, že \(\Phi\) nastaje tvorjenjem polarov (sravni formulu (2) na citovanom městu); polny razklad v rědy znovu dobiva se iteracijeju.

Tvorjenje normalnyh form \(\Phi\) zato v obah slučajah podređuje se slědujučemu problemu. Ako v \(F(z)\) neizvěstne \(z_i\) zaměnjajut se polynomami \(f_i(\xi)\) od parametrov \(\xi\), tako že predstavjenje \(F(f_i(\xi))\) uže ne jest jednoznačno, togda jednoznačnosť treba ustanoviti tym, že medžu koeficientami važe te same linearne relacije, kako medžu produktami stepenev \(f_i(\xi)\), ktore nastupajut v \(F\). Tako imamo
\begin{equation}
\begin{aligned}
        F(f_i(\xi))&=\Phi(f_i(\xi))&&[\text{identično v }\xi],\quad\text{ili}\notag\\
        F(z)&\equiv\Phi(z)&&\modu{M},
\end{aligned}
\tag{3}
\end{equation}
gdě \(M\) označaje modul vseh identičnyh relacij medžu \(f_i(\xi)\).

Za tutu normalnu formu \(\Phi(z)\) E. Fišer v \S~11, I svojej citovanoj raboty dal eksplicitny razklad do čislovyh koeficientov, imenno pomoću linearnyh operacij. Jednovrěmenno v uvodu III i v \S~6 II dal on za razliku \(F-\Phi\), naležeču k \(M\), pri predpoloženju znanja bazisa modula, eksplicitny izraz v tom samom smyslu; i v povezivanju tyh razkladov v \S~11 II zaměnil formulu [3] eksplicitnym predstavjenjem. Pri tom ješče otvoreny problem čislovyh koeficientov po \S~12 toj samej raboty jest identičny s iskanjem vlastnyh vrědnostij i vlastnyh funkcij linearnyh operacij (formula 82).

Teper v punktu 1 specializujem Fišerove formuly na slučaj razklada v rědy po polarah i dobivam iz njih, čerez iteraciju, razklad v rědy. Operacija za opreděljenje \(\Phi\) pri tom prehodi v opreděljene polarny procesy, čto jest mnogo točnejše než gore podana, dotolě znana formulacija. V operaciji za opreděljenje \(F-\Phi\) směšano nastupajut množenje determinantem \(\Delta\) i \(\Omega\)-proces. Da by se došlo k obyčnoj, neeksplicitnoj formě, treba ješče zamětiti, že \(\Omega\Delta\) možno izraziti polarnymi procesami.

Na tutom městě treba dodati ješče jedno popravjenje. V citovanoj rabotě (s. 101, ręd 6 od vrha) slučajno prijela jesm identično prětvorjenje
\[
        \Omega(\Delta\psi)=c_1\psi+c_2\Delta\cdot\Omega(\psi),
\]
ktoro važi samo v binarnoj oblasti, kako obče važeče, i tym sčinila prěhod od fundamentalnoj formuly (2), ktora daje izraz za normalnu formu \(\Phi\) v [2], k razkladu v rědy (3). Tamošnje razsmatranja zato dajut samo formulu (2), i kako specialny slučaj iz njej redukcijny teorem (1); ale, poněže manjka eksplicitny izraz za razliku \(F-\Phi\), ne dostajut za osnovanje razklada v rědy.

V punktu 2 dobivam ravno tako specializacijeju normalnu formu \(\Phi(t_{ik})\) iz [1]; občeje, normalnu formu \(\Phi(t_i,t_{ik},t_{ikl},\ldots)\), čto vodi k znanym eksplicitnym formulam. Ale poněže postavjenje polnogo razklada v rědy potrěbuje znanje bazisa modula \(M\), tude jest poželatelnějši invariantno-teoretičny put (sravni navedenu literaturu), ktory ne potrěbuje to znanje, ale naprotiv, pomoću razklada v rědy daje bazis modula.

1. Po Fišeru, \S~11 (formula (19) i slědujuče rědy), normalna forma \(\Phi(z)\) iz [3], tam \(N(\zeta)\), dana jest:
\begin{equation}
        \Phi(z)=(a_1S+a_2S^2+\cdots+a_rS^r)F(z),
\tag{4}
\end{equation}
gdě konstanty \(a_1,\ldots,a_r\) naležet k oblasti koeficientov \(f(\xi)\) i zaviset samo od stepena \(F(z)\); operator \(S\) opreděljen jest čerez
\begin{equation}
\begin{gathered}
        S=AB;\qquad BF(z)=F(f(\xi));\\
        AF(f(\xi))=\left\{\sum z_i f_i\!\left(\frac{\p}{\p\xi}\right)\right\}^{p}F(f(\xi))\footnotemark.
\end{gathered}
\tag{5}
\end{equation}
\footnotetext{Sravni formulu (75). V slučaju [3] suma \(\alpha\) proběgaje vse produkty stepenev \(p\)-go razměra, i zato (75) specializuje se do [5].}
Ako v [2] \(F\) pojmati kako formu \(p\)-go razměra samo v \(z\), togda \(f_i(\xi)\) specializuje se do \(\lambda_1x_i+\lambda_2y_i\); zato dobiva se:
\[
\begin{gathered}
        BF=F(x,y,\lambda_1x+\lambda_2y),\\
        SF=\left\{\sum z_i\left(\frac{\p}{\p\lambda_1}\frac{\p}{\p x_i}+
        \frac{\p}{\p\lambda_2}\frac{\p}{\p y_i}\right)\right\}^{p}
        F(x,y,\lambda_1x+\lambda_2y),
\end{gathered}
\]
ili, zapisano v polarnyh procesah (sravni cit. město II, 2 i 3; takože za skraćeno označenje):
\begin{equation}
        SF(x,y,z)=\frac1{p!}\left[z\left(\frac{\p}{\p\lambda_1}\frac{\p}{\p x}+
        \frac{\p}{\p\lambda_2}\frac{\p}{\p y}\right)\right]^p
        \left[(\lambda_1x+\lambda_2y)\frac{\p}{\p z}\right]^p F(x,y,z).
\tag{6}
\end{equation}
{\itshape Formulami} [4] {\itshape i} [6] {\itshape jest dany potrěbny eksplicitny izraz za} \(\Phi(x,y,z)\); \(\Phi(x,y,z)\) jednovrěmenno podređuje se obče formě spomenutoj na početku (\(\sum PH\) v (2), cit. město), polarom izrazov, ktore zaviset samo od 2 rędov i dobivajut se iz \(F\) polarnymi procesami. Ješče neopreděljene koeficienty \(a_i\) sut racionalne čisla, poněže tude \(f_i(\xi)\) imajut cělo-racionalne čislove koeficienty. Ako zapisati [2] v formě
\begin{equation}
        F=\Phi+\Delta F_1,
\tag{2a}
\end{equation}
togda specializacijeju formuly, podanoj u Fišera v uvodu III i v \S~6 II (s. 41), dobiva se
\begin{equation}
        F_1=(c_1\Omega+c_2\Omega\Delta\Omega+\cdots+c_i\Omega\Delta\Omega\cdots\Delta\Omega)F\footnotemark,
        \qquad \left[\Omega=\left(\frac{\p}{\p x}\frac{\p}{\p y}\frac{\p}{\p z}\right)\right],
\tag{7}
\end{equation}
\footnotetext{Tutu specializovanu formulu upotrěbljaje Kerim v svojej skoro izhodečej erlangenskoj disertaciji v priměnjenju na «formy inercije».}
gdě \(c_i\) znovu sut racionalne čisla, zaviseče samo od stepena \(F\).

Iteracija vodi k
\[
        F_1=\Phi_1+\Delta F_2;\quad F_2=\Phi_2+\Delta F_3;\quad\ldots;\quad
        F_i=\Phi_i+\Delta F_{i+1};\quad\ldots,
\]
gdě vsaky raz \(\Phi_i\) sut normalne formy odpovědajuče k \(F_i\). Poněže \(F_i\) imaje v \(z\) stepen \((p-i)\), po [4] i [6] dobiva se:
\begin{equation}
\begin{gathered}
        \Phi_i=(a_{i1}S_i+a_{i2}S_i^2+\cdots)F_i,\\
        S_iF_i=\frac1{(p-i)!}
        \left[z\left(\frac{\p}{\p\lambda_1}\frac{\p}{\p x}+
        \frac{\p}{\p\lambda_2}\frac{\p}{\p y}\right)\right]^{p-i}
        \left[(\lambda_1x+\lambda_2y)\frac{\p}{\p z}\right]^{p-i}F_i.
\end{gathered}
\tag{8}
\end{equation}
Formuli [7] odpovědaje
\[
        F_i=(\alpha_{i1}\Omega+\alpha_{i2}\Omega\Delta\Omega+\cdots)F_{i-1},
        \quad\text{i čerez iteraciju}
\]
\begin{equation}
        F_i=(\beta_1\Omega^i+\beta_2\Omega^i\Delta\Omega+\cdots
        +\beta_\rho\Omega^{k_1}\Delta\Omega^{k_2}\cdots\Delta\Omega^{k_\lambda}+\cdots)F,
        \quad(k_1+k_2+\cdots+k_\lambda-\lambda+1=i).
\tag{9}
\end{equation}
{\itshape Pomoću} [8] {\itshape i} [9] {\itshape v razkladu}
\[
        F=\Phi+\Delta\Phi_1+\Delta^2\Phi_2+\cdots+\Delta^i\Phi_i+\cdots+\Delta^p\Phi_p
\]
{\itshape dostignuto jest eksplicitno predstavjenje v naznačenom smyslu.}

Ako na [7] priměniti identično prětvorjenje \(\Omega\Delta F=PF\), gdě \(P\) označaje polarny procesy,\footnote{Sravni, na primer, F. Mertens, \emph{Über ganze Funktionen von \(m\) Systemen von je \(n\) Unbestimmten}, Monatsh. f. Math. u. Phys. 4, formula (5).} togda, zaradi zaměnljivosti tyh procesov s \(\Omega\)-procesom, dobiva se znana formula
\[
        F_1=P\Omega F;\quad\ldots;\quad F_i=P\Omega^iF;
\]
i iz togo v povezku s [8] ili s razsmatranjami na citovanom městu slěduje znana forma razklada v rědy, kako na primer v (3) na citovanom městu. Ako imajemo \(n\) rędov po \(n\) proměnnyh, \(x^{(1)},x^{(2)},\ldots,x^{(n-1)},z\), togda odpovědajuče
\[
        f_i(\xi)=\lambda_1x_i^{(1)}+\lambda_2x_i^{(2)}+\cdots+\lambda_{n-1}x_i^{(n-1)};
\]
i jednako treba opreděliti \(\Delta\) i \(\Omega\) za \(n\) rędov.

2. Za normalnu formu \(F(t_i,t_{ik},t_{ikl},\ldots)\) funkcije \(f_i(\xi),f_{ik}(\xi),\ldots\) dane sut jedno-, dvo-, do \((n-1)\)-rędnyh determinantov
\[
        \xi_i^{(1)},\quad (\xi^{(1)}\xi^{(2)})_{ik},\quad\ldots,\quad
        (\xi^{(1)}\xi^{(2)}\cdots\xi^{(n-1)})_{i_1i_2\cdots i_{n-1}}.
\]
Zato kako specializacija [5] dobiva se:
\[
\begin{gathered}
        S=AB;\qquad
        BF=F\bigl(\xi_i^{(1)},(\xi^{(1)}\xi^{(2)})_{ik},\ldots\bigr),\\
        ABF=\left(\sum t_i\frac{\p}{\p\xi_i^{(1)}}\right)^{\lambda_1}
        \left(\sum t_{ik}\left(\frac{\p}{\p\xi^{(1)}}\frac{\p}{\p\xi^{(2)}}\right)_{ik}\right)^{\lambda_2}
        \cdots F\bigl(\xi_i^{(1)},(\xi^{(1)}\xi^{(2)})_{ik},\ldots\bigr),
\end{gathered}
\]
gdě \(\lambda_1,\lambda_2,\ldots\) označajut stepeni v \(t_i,t_{ik},\ldots\).

Iz invariantno-teoretičnyh razsmatranj, imenno že vse invarianty \(F\), ktore sut linearne v koeficientah, možno linearno sostaviti iz \(\Phi\) i iz form iz \(M\), zato osobno takože \(SF\), ktoro ne naleži k \(M\),\footnote{To slěduje na primer iz \S~6 i \S~7, 2 moje citovanoj raboty \emph{Zur Invariantentheorie der Formen von \(n\) Variabeln}.} tude slěduje kongruencija \(\Phi\equiv a_1SF\modu{M}\); i zaradi vlastnosti normalnoj formy iz togo vměsto [4] dobiva se prosta relacija
\begin{equation}
        \Phi=a_1SF=a_1S\Phi.
\tag{10}
\end{equation}
Drugo ravnanje važi, poněže priměnjenje \(S\) vodi vsaku formu iz \(M\) k nuli. Formula [10] pokazuje, že sama normalna forma \(\Phi\) jednovrěmenno jest vlastna funkcija; i po gornjem invariantnom razsmatranju ne obstavajut druge vlastne funkcije, ktore by jednovrěmenno byly invariantami \(F\) (sravni k [10] Derujts, formuly (10) i (10')).

\begin{center}
\textbf{Popravjenja.}
\end{center}

1. V rabotě \emph{Die allgemeinsten Bereiche aus ganzen transzendenten Zahlen} (\emph{Math. Ann.} 77), na s. 121, ręd 8 od vrha, ręd 9 od vrha i ręd 11 od dola, polje \((H,Y)\) treba zaměniti na «polje \((H,Y,Y',\ldots,Y^{(\lambda_0-1)})\), gdě \(Y',\ldots,Y^{(\lambda_0-1)}\) označajut spregnute od \(Y\) vzhodno k \(H\)». Odpovědajuče, na s. 120, ręd 5 od vrha, i s. 121, ręd 9 od vrha, \((H,y)\) treba zaměniti na \((H,y,y',\ldots,y^{(\lambda_0-1)})\).

2. V rabotě \emph{Gleichungen mit vorgeschriebener Gruppe} (\emph{Math. Ann.} 78), pri formulaciji Hilbertovogo teorema o nerazložimosti, s. 222, ręd 3 od dola, vměsto «čislovo polje \(\Omega\)» točnejše treba kazati: «konečno algebraično čislovo polje \(\Omega\)», na čto mi ukazuje gospodin Ostrovski. V samoj stvari, na primer vzhodno k polju vsih algebraičnyh čisel, grupa vsakogo uravnenja s algebraičnymi čislovymi koeficientami jest identična grupa. Odpovědajuče, v uvodu pod «libovoljnoju oblastju racionalnosti» treba razuměti konečno algebraično čislovo polje, k ktoromu ješče možno adjungovati konečno mnogo parametrov, ili jednu iz tako nastajučih medžuoblastij.

\begin{center}
(Prijeto v septembru 1919.)
\end{center}

\end{document}
